\section{The sharp lower constant at every length}\label{sec:lower}
First use the finite comparison with
\[
 r=\lfloor2n/\log n\rfloor,\qquad q=n-r,\qquad M=r(r+1)/2.
\]
Writing $L=\log n$, these parameters satisfy
\begin{equation}\label{eq:optimparameters}
 r=\frac{2n}{L}+O(1),\qquad M\sim\frac{2n^2}{L^2},\qquad
 Q\sim\frac{4n}{L^2}.
\end{equation}
The three fiber costs are subexponential. More explicitly,
\begin{align}
 \log\left(2^{r+1}(Q+1)\binom{n+Q}{r}\right)
 &\le(r+1)\log2+\log(Q+1)
      +r\log\frac{e(n+Q)}r\notag\\
 &=O(n\log L/L).\label{eq:lowercost}
\end{align}
By Lemma~\ref{lem:triangular},
\begin{equation}\label{eq:singlelower}
 \log T(n,r)=n(L-2\log L+C_+)+O(n\log L/L).
\end{equation}
Equations \eqref{eq:finitelower}--\eqref{eq:singlelower} give the desired lower constant at lengths $n+Q$, but those lengths alone are not enough for an all-length theorem.

Fix a sufficiently large target $N$, and choose instead
\begin{equation}\label{eq:alllength}
 n=N-\left\lceil\frac{8N}{(\log N)^2}\right\rceil,
 \qquad r=\lfloor2n/\log n\rfloor,
\end{equation}
with $Q$ as in \eqref{eq:Q}. Then $n/N\to1$, $Q\sim4N/(\log N)^2$, and consequently $n+Q\le N$ eventually. Apply Proposition~\ref{prop:repair}, then append additional fresh records to reach exactly $N$. This extra padding is injective because its starting length is fixed.

For $f(x)=x(\log x-2\log\log x+C_+)$, the explicit derivative is
\[
 f'(x)=\log x-2\log\log x+C_++1-\frac2{\log x}=O(\log x).
\]
The shift in \eqref{eq:alllength} is $O(N/(\log N)^2)$, so
\[
 f(n)=f(N)+O(N/\log N).
\]
All the losses in \eqref{eq:lowercost} are $O(N\log\log N/\log N)$. It follows that
\begin{equation}\label{eq:lowerfinal}
 \log e_N\ge N(\log N-2\log\log N+C_+)
                  -O(N\log\log N/\log N).
\end{equation}
Since $e_N\le a_N$ and both satisfy the upper bound \eqref{eq:upperfinal}, this proves \eqref{eq:main} for $a,e$. In particular, the stronger lower construction lands entirely in the multiplicity-at-most-two subclass $\Av(000,100)$.

Combining \eqref{eq:main} with \eqref{eq:blog} also yields
\[
 \log a_n-\log b_n=o(n),\qquad \log e_n-\log b_n=o(n).
\]
This means $a_n/b_n=\exp(o(n))$ and $e_n/b_n=\exp(o(n))$, often called subexponential equivalence. It is strictly weaker than either ratio tending to one, and gives no relative equivalent for $a_n$ or $e_n$.
